% Source: https://artofproblemsolving.com/community/c3h1927516_question_about_multiplicities
% Stable identity: aops-topic-1927516-q01
\textbf{(a)} Since $x^2+36=(x-6i)(x+6i)$,
\[
f(x)=(x-2)(x-1)(x-6i)(x+6i).
\]
The roots are $\boxed{2,\ 1,\ 6i,\ -6i}$, each of multiplicity $1$. (The multiplicities add up to the degree, $4$.)

\textbf{(b)} No. Multiplicity is a property of a root, not of a factor. The quadratic factor $x^2+36$ contributes two different roots, $6i$ and $-6i$, each of multiplicity $1$. The multiplicity of a root $r$ equals the exponent of $(x-r)$ only after the polynomial is factored completely into linear factors and equal factors are collected.

\textbf{(c)} $g(x)=(x-2)^3(x-6i)^2(x+6i)^2$, so $2$ has multiplicity $3$, and $6i$ and $-6i$ each have multiplicity $2$ (total $7$, the degree of $g$).
